% ===== PRÉAMBULE CANONIQUE — à coller VERBATIM en tête de CHAQUE .tex =====
\documentclass[11pt,a4paper]{article}
\usepackage[utf8]{inputenc}\usepackage[T1]{fontenc}\usepackage{lmodern}
\usepackage[french]{babel}
\usepackage{amsmath,amssymb,mathtools}\usepackage{icomma}
\usepackage{siunitx}\sisetup{locale=FR}  % \qty{}{}, \si{}, \ang{}
\usepackage{xcolor}\usepackage{colortbl}
\usepackage{tikz}\usepackage{pgfplots}\pgfplotsset{compat=1.18}
\usepackage[siunitx,europeanresistors]{circuitikz} % circuits (élec.)
\usepackage[most]{tcolorbox}\usepackage{enumitem}\usepackage{array,booktabs}
\usepackage[a4paper,margin=2.2cm]{geometry}
\usetikzlibrary{arrows.meta,calc,angles,quotes,positioning,patterns,%
  backgrounds,shapes.geometric,decorations.pathmorphing,decorations.markings,intersections}
\definecolor{primary}{RGB}{222,49,99}\definecolor{primarydark}{RGB}{150,22,55}
\definecolor{defcol}{RGB}{0,114,178}\definecolor{methcol}{RGB}{52,92,148}
\definecolor{excol}{RGB}{0,135,90}\definecolor{attncol}{RGB}{200,40,55}
\tcbset{boxbase/.style={enhanced,breakable,boxrule=0.9pt,arc=2.5pt,
  left=3mm,right=3mm,top=2.3mm,bottom=2.3mm,fonttitle=\bfseries,coltitle=white}}
% --- boîtes TUTORIEL (noms EXACTS, ne pas renommer) ---
\newtcolorbox{cle}[1][]{boxbase,colback=defcol!6,colframe=defcol,
  title={\textsc{À retenir}\ifx\relax#1\relax\relax\else\textmd{ : \itshape #1}\fi}}
\newtcolorbox{methode}[1][]{boxbase,colback=methcol!7,colframe=methcol,sharp corners,
  title={\textsc{Méthode}\ifx\relax#1\relax\relax\else\textmd{ --- \itshape #1}\fi}}
\newtcolorbox{exemplebox}[1][]{boxbase,colback=excol!5,colframe=excol,
  title={\textsc{Exemple}\ifx\relax#1\relax\relax\else\textmd{ : \itshape #1}\fi}}
\newtcolorbox{piege}[1][]{boxbase,colback=attncol!6,colframe=attncol,
  title={\textsc{Piège}\ifx\relax#1\relax\relax\else\textmd{ : \itshape #1}\fi}}
\newtcolorbox{remarque}[1][]{boxbase,colback=black!4,colframe=black!45,
  title={\textsc{Remarque}\ifx\relax#1\relax\relax\else\textmd{ : \itshape #1}\fi}}
% --- boîtes EXERCICE / CORRIGÉ (noms EXACTS, auto-comptées) ---
\newtcolorbox[auto counter]{exo}[1][]{boxbase,colback=primary!4,colframe=primary,
  title={\textsc{Exercice}~\thetcbcounter\ifx\relax#1\relax\relax\else\textmd{ --- \itshape #1}\fi}}
\newtcolorbox[auto counter]{corrige}[1][]{boxbase,colback=excol!4,colframe=excol,
  title={\textsc{Corrigé}~\thetcbcounter\ifx\relax#1\relax\relax\else\textmd{ --- \itshape #1}\fi}}
\newcommand{\vv}[1]{\vec{#1}}\newcommand{\ds}{\displaystyle}
\newcommand{\R}{\mathbb{R}}\newcommand{\N}{\mathbb{N}}\newcommand{\Z}{\mathbb{Z}}
% (un \usepackage{fancyhdr}+en-tête est possible ; il est ignoré côté web)
% =========================================================================
\begin{document}
\sisetup{per-mode=symbol}

\begin{corrige}[champ d'une charge ponctuelle]
\begin{enumerate}[label=\alph*)]
  \item \[ E = K\,\frac{|Q|}{r^2} = \num{9e9}\times\frac{\num{8e-9}}{(0{,}20)^2}
             = \frac{72}{0{,}04} = \SI{1800}{\newton\per\coulomb}. \]
  \item La charge est \emph{positive} : $\vv{E}$ \textbf{fuit} la charge, il pointe donc à l'opposé
        de $Q$ (vers l'extérieur).
\end{enumerate}
\end{corrige}

\begin{corrige}[de la force au champ]
\begin{enumerate}[label=\alph*)]
  \item \[ E = \frac{F}{q} = \frac{\num{8e-6}}{\num{4e-9}} = \SI{2000}{\newton\per\coulomb}. \]
  \item Le champ est le \emph{même} : $E = \SI{2000}{\newton\per\coulomb}$. Le champ \textbf{ne
        dépend pas} de la charge test (avec $q'$, la force doublerait, mais $F/q'$ redonne la même
        valeur).
\end{enumerate}
\end{corrige}

\begin{corrige}[la force dans un champ]
\begin{enumerate}[label=\alph*)]
  \item $F = |q|\,E = \num{2e-6}\times 500 = \SI{1e-3}{\newton}$. Comme $q>0$, $\vv{F}$ a le
        \emph{même} sens que $\vv{E}$ : vers la \textbf{droite}.
  \item Même intensité $F = \SI{1e-3}{\newton}$, mais $q<0$ : $\vv{F}$ est de sens \emph{opposé} à
        $\vv{E}$, vers la \textbf{gauche}.
\end{enumerate}
\end{corrige}

\begin{corrige}[orienter le champ]
\begin{enumerate}[label=\alph*)]
  \item $\vv{E}$ \emph{fuit} la charge positive : en $M_1$ (à droite du $+$), il pointe vers la
        \textbf{droite} (en s'éloignant du $+$).
  \item $\vv{E}$ \emph{pointe vers} la charge négative : en $M_2$ (à gauche du $-$), il pointe vers
        la \textbf{droite} (vers le $-$).
\end{enumerate}
\end{corrige}

\begin{corrige}[champ nul au milieu]
\begin{enumerate}[label=\alph*)]
  \item $C$ est à la même distance $a$ de $A$ et de $B$, et les charges sont égales :
        \[ E_A = E_B = K\,\frac{q}{a^2}\quad\text{(mêmes intensités)}. \]
  \item Les deux champs \emph{fuient} leur charge : $\vv{E}_A$ pointe vers la droite, $\vv{E}_B$
        vers la gauche. Ils sont opposés et de même norme :
        \[ \vv{E}_{\text{tot}}(C) = \vv{E}_A + \vv{E}_B = \vv{0}. \]
\end{enumerate}
\end{corrige}

\end{document}
